\documentclass[11pt]{amsart}
\usepackage[T1]{fontenc}
\usepackage{amsmath,amssymb,amsthm}
\usepackage[colorlinks=true,linkcolor=blue,citecolor=blue,urlcolor=blue]{hyperref}

\newtheorem{theorem}{Theorem}
\newtheorem{lemma}{Lemma}
\newtheorem{proposition}{Proposition}
\newtheorem{corollary}{Corollary}

\newcommand{\CC}{\mathcal{CC}}
\newcommand{\LLL}{\mathcal{L}}
\newcommand{\Lmsc}{\mathcal{L}_{msc}}
\newcommand{\tensorpi}{\widehat{\otimes}_{\pi}}

\begin{document}

\title{A Conditional Reduction for the Hyper-Ideal Question}
\author{Run fa\_banach\_001}
\date{2026-06-16}
\maketitle

\section*{Status}

This packet proves a conditional reduction. It does not prove the external
Schur tensor-product dependency. Conditional status:
\[
  \texttt{conditional\_reduction\_likely\_valid}.
\]

\section*{Source Question}

Botelho--Torres, \emph{Hyper-ideals of multilinear operators},
arXiv:1504.00520, define the hyper-ideal \(\Lmsc\) of multilinearly
sequentially continuous multilinear operators. They prove
\[
  \CC\circ\LLL \subseteq \Lmsc,
\]
where \(\CC\) is the closed ideal of completely continuous linear operators
and \(\CC\circ\LLL\) is the composition hyper-ideal.

Their Open Question 2 asks whether this inclusion is proper:
\[
  \CC\circ\LLL \subsetneqq \Lmsc.
\]
The same discussion relates this question to the long-standing problem of
whether the Schur property is preserved by projective tensor products.

\section*{Conditional Dependencies}

The only external dependency is the following unproved assumption.

\medskip
\noindent
\textbf{Schur tensor counterexample hypothesis.}
There exist Banach spaces \(E_1,\ldots,E_n\), \(n\geq 2\), each with the
Schur property, such that
\[
  Z=E_1\tensorpi\cdots\tensorpi E_n
\]
does not have the Schur property.

\medskip
Rodriguez, \emph{On weak compactness in projective tensor products},
arXiv:2206.08651, still records Schur preservation under projective tensor
products as open. Thus this packet is not a full answer to Open Question 2;
it proves that any negative answer to that tensor-product problem gives the
desired proper inclusion.

\section*{Proof Intuition}

The composition ideal \(\CC\circ\LLL\) can be detected on the linearization of
a multilinear map: an \(n\)-linear operator belongs to \(\CC\circ\LLL\) exactly
when its associated linear operator on the projective tensor product is
completely continuous. The canonical tensor map has identity as its
linearization, so it fails to be in \(\CC\circ\LLL\) precisely when the tensor
product lacks the Schur property.

The remaining issue is whether the canonical tensor map lies in \(\Lmsc\).
For Schur factors, weak convergence of elementary tensors is already norm
convergence. The zero-limit case is forced by a subsequence/functionals
argument using the Schur property in each factor; the nonzero case is reduced
to the zero-limit case by splitting the factors into rank-one limit directions
and complementary kernels.

\section*{Elementary Tensors over Schur Factors}

\begin{lemma}\label{lem:zero}
Let \(E_1,\ldots,E_n\) have the Schur property. If
\[
  z_j=x_j^1\otimes\cdots\otimes x_j^n
\]
converges weakly to \(0\) in \(E_1\tensorpi\cdots\tensorpi E_n\), then
\[
  \|x_j^1\|\cdots\|x_j^n\|\longrightarrow 0.
\]
\end{lemma}

\begin{proof}
The sequence \((z_j)\) is weakly convergent, hence bounded. Suppose, after
passing to a subsequence, that
\[
  \|x_j^1\|\cdots\|x_j^n\|\geq \delta>0
\]
for all \(j\). Terms with a zero factor can be discarded, so put
\[
  u_j^k=x_j^k/\|x_j^k\| \qquad (1\leq k\leq n).
\]
For every continuous \(n\)-linear form \(B\),
\[
  B(u_j^1,\ldots,u_j^n)
  =
  \frac{B(x_j^1,\ldots,x_j^n)}
       {\|x_j^1\|\cdots\|x_j^n\|}
  \longrightarrow 0.
\]

Since \(E_1\) has the Schur property and \(\|u_j^1\|=1\), the sequence
\((u_j^1)\) is not weakly null. Hence some subsequence and some
\(\varphi_1\in E_1^*\) satisfy \(|\varphi_1(u_j^1)|\geq \epsilon_1>0\) on
that subsequence. Repeating this argument successively in \(E_2,\ldots,E_n\),
we obtain a common subsequence and functionals \(\varphi_k\in E_k^*\) such
that \(|\varphi_k(u_j^k)|\geq \epsilon_k>0\) for every \(k\) along the common
subsequence. The multilinear form
\[
  B(y_1,\ldots,y_n)=\varphi_1(y_1)\cdots\varphi_n(y_n)
\]
then does not tend to \(0\) on that subsequence, a contradiction.
\end{proof}

\begin{lemma}\label{lem:pure-schur}
Let \(E_1,\ldots,E_n\) have the Schur property. If elementary tensors
\[
  z_j=x_j^1\otimes\cdots\otimes x_j^n
\]
converge weakly to an elementary tensor
\[
  z=x^1\otimes\cdots\otimes x^n
\]
in \(E_1\tensorpi\cdots\tensorpi E_n\), then \(z_j\to z\) in projective norm.
\end{lemma}

\begin{proof}
If \(z=0\), this is Lemma~\ref{lem:zero}. Suppose \(z\neq 0\), so every
\(x^k\neq 0\). Choose \(\varphi_k\in E_k^*\) with \(\varphi_k(x^k)=1\).
Let \(P_k(y)=\varphi_k(y)x^k\) and \(Q_k=I_{E_k}-P_k\). For each subset
\(S\subseteq\{1,\ldots,n\}\), let
\[
  R_S=T_1^S\tensorpi\cdots\tensorpi T_n^S,
  \qquad
  T_k^S=
  \begin{cases}
    Q_k, & k\in S,\\
    P_k, & k\notin S.
  \end{cases}
\]
The operators \(R_S\) are bounded and weak-to-weak continuous on the
projective tensor product. Also \(R_{\varnothing}z=z\) and \(R_S z=0\) for
\(S\neq\varnothing\).

Since \(z_j\to z\) weakly, \(R_Sz_j\to 0\) weakly for every nonempty \(S\).
Each \(R_Sz_j\) is an elementary tensor in a projective tensor product of
closed subspaces of the \(E_k\)'s, hence again in Schur spaces. By
Lemma~\ref{lem:zero}, \(\|R_Sz_j\|\to 0\) for every nonempty \(S\).

For \(S=\varnothing\),
\[
  R_{\varnothing}z_j
  =
  \left(\prod_{k=1}^n \varphi_k(x_j^k)\right)
  x^1\otimes\cdots\otimes x^n
  \longrightarrow z
\]
in norm, because applying the scalar \(n\)-linear form
\((y_1,\ldots,y_n)\mapsto\prod_k\varphi_k(y_k)\) gives convergence of the
coefficient to \(1\). The finite decomposition
\[
  z_j=\sum_{S\subseteq\{1,\ldots,n\}}R_Sz_j
\]
therefore gives \(z_j\to z\) in projective norm.
\end{proof}

\section*{The Conditional Separation}

\begin{proposition}\label{prop:linearization}
For \(A\in\LLL(E_1,\ldots,E_n;F)\), let
\(\widetilde A:E_1\tensorpi\cdots\tensorpi E_n\to F\) be its linearization.
Then
\[
  A\in \CC\circ\LLL(E_1,\ldots,E_n;F)
\]
if and only if \(\widetilde A\) is completely continuous.
\end{proposition}

\begin{proof}
If \(A=u\circ B\) with \(u:G\to F\) completely continuous and
\(B:E_1\times\cdots\times E_n\to G\) multilinear, then
\(\widetilde A=u\circ\widetilde B\). Since \(\widetilde B\) is weak-to-weak
continuous and \(u\) sends weakly convergent sequences to norm-convergent
sequences, \(\widetilde A\) is completely continuous.

Conversely, if \(\widetilde A\) is completely continuous, then
\[
  A=\widetilde A\circ\sigma,
  \qquad
  \sigma(x_1,\ldots,x_n)=x_1\otimes\cdots\otimes x_n,
\]
is a \(\CC\circ\LLL\) factorization.
\end{proof}

\begin{theorem}\label{thm:conditional}
Assume the Schur tensor counterexample hypothesis. Then
\[
  \CC\circ\LLL \subsetneqq \Lmsc.
\]
\end{theorem}

\begin{proof}
Let \(E_1,\ldots,E_n\) be Schur spaces such that
\[
  Z=E_1\tensorpi\cdots\tensorpi E_n
\]
does not have the Schur property. Consider the canonical tensor map
\[
  \sigma:E_1\times\cdots\times E_n\to Z,\qquad
  \sigma(x_1,\ldots,x_n)=x_1\otimes\cdots\otimes x_n.
\]

If a sequence of tuples converges multilinearly, then the corresponding
elementary tensors converge weakly in \(Z\). By Lemma~\ref{lem:pure-schur},
they converge in projective norm. Hence \(\sigma\in\Lmsc(E_1,\ldots,E_n;Z)\).

The linearization of \(\sigma\) is the identity operator \(I_Z\). Since \(Z\)
lacks the Schur property, \(I_Z\) is not completely continuous. By
Proposition~\ref{prop:linearization}, \(\sigma\notin
\CC\circ\LLL(E_1,\ldots,E_n;Z)\). Thus the inclusion is proper.
\end{proof}

\section*{Verification Notes}

The proof above is complete modulo the explicitly stated Schur tensor
counterexample hypothesis. No computational evidence is used. The most
important review point is Lemma~\ref{lem:pure-schur}, especially the reduction
from a nonzero elementary limit to the zero-limit lemma by tensoring the
rank-one projections \(P_k\) and \(Q_k\).

\section*{Search Bounds}

The run indexes were searched for arXiv:1504.00520, hyper-ideals,
multilinearly sequentially continuous operators, \(\CC\circ\LLL\), and
projective tensor Schur phrases. No prior packet for this conditional
reduction was found. Local parsed arXiv sources and a bounded web search on
2026-06-16 did not identify a later exact solution of the Schur
tensor-product preservation problem; Rodriguez arXiv:2206.08651 still records
it as open in 2022.

\begin{thebibliography}{9}
\bibitem{BT}
G. Botelho and E. R. Torres,
\emph{Hyper-ideals of multilinear operators},
arXiv:1504.00520.

\bibitem{Rodriguez}
J. Rodriguez,
\emph{On weak compactness in projective tensor products},
arXiv:2206.08651.
\end{thebibliography}

\end{document}
